% 第5章图：作用域{1,2,3}与{2,3,4}的因子图、原始关系图与超图。
\begin{tikzpicture}[x=1mm,y=1mm,>=Stealth,
  var/.style={circle,draw=FigureStroke,fill=FigureBlueFill,line width=.8pt,minimum size=6.5mm,inner sep=0pt},
  fac/.style={rectangle,draw=FigureStroke,fill=FigureGreenFill,line width=.8pt,minimum size=5.5mm,inner sep=1pt},
  e/.style={FigureStroke,line width=.8pt},
  ttl/.style={font=\figurefont\footnotesize,text=FigureInk}]
% (a) factor graph: variables in a row, factors above
\begin{scope}
 \foreach \i in {1,2,3,4}{\node[var] (a\i) at ({(\i-1)*9},0) {$\i$};}
 \node[fac] (fA) at (9,16) {$A$};
 \node[fac] (fB) at (18,16) {$B$};
 \draw[e] (fA)--(a1) (fA)--(a2) (fA)--(a3) (fB)--(a2) (fB)--(a3) (fB)--(a4);
 \node[ttl] at (13.5,-15) {(a) 因子图};
\end{scope}
% (b) primal graph
\begin{scope}[xshift=38mm,yshift=-1mm]
 \node[var] (b1) at (0,9) {$1$};
 \node[var] (b2) at (12,18) {$2$};
 \node[var] (b3) at (12,0) {$3$};
 \node[var] (b4) at (24,9) {$4$};
 \draw[e] (b1)--(b2) (b1)--(b3) (b2)--(b3) (b2)--(b4) (b3)--(b4);
 \draw[FigureMuted,dashed,line width=.6pt] (b1)--(b4);
 \node[font=\FigureNoteFont\fontsize{7.5}{9}\selectfont,text=FigureMuted,fill=white,inner sep=.8pt] at (7.6,9) {无边};
 \node[ttl] at (12,-14) {(b) 原始关系图};
\end{scope}
% (c) hypergraph
\begin{scope}[xshift=75mm,yshift=-1mm]
 \draw[FigureBlue,line width=.8pt,rounded corners=5mm] (-8,9)--(17,28)--(17,-10)--cycle;
 \draw[FigureGreen,line width=.8pt,dashed,rounded corners=5mm] (7,28)--(32,9)--(7,-10)--cycle;
 \node[var] at (0,9) {$1$};
 \node[var] at (12,18) {$2$};
 \node[var] at (12,0) {$3$};
 \node[var] at (24,9) {$4$};
 \node[ttl] at (12,-14) {(c) 超图};
\end{scope}
\end{tikzpicture}
